%% kapa: exam=1 round=21 year=2010 period=2 form=A subject=appraisal_law subject_ko=부동산관계법규 seq=35 number=75
%% source: Q-Net 감정평가사 기출문제 (cst00309) artlSeq=5024181 file=extracted/1008/5024181/2010년 제21회 감정평가사 1차시험 기출문제/2교시.hwp
%% conversion: kordoc 4.24.51 -> tools/gen1.py (auto); review=proofread-gpt-sol; flags=none
\begin{question}{75}
측량․수로조사 및 지적에 관한 법령상 토지소유자의 합병 신청이 제한되는 경우가 아닌 것은?
\begin{choices}
  \item 합병하려는 토지의 지번부여지역이 서로 다른 경우
  \item 합병하려는 토지에 지상권의 등기가 있는 경우
  \item 합병하려는 토지의 지목이 서로 다른 경우
  \item 합병하려는 토지의 소유자가 서로 다른 경우
  \item 합병하려는 토지의 지적도 및 임야도의 축척이 서로 다른 경우
\end{choices}
\end{question}
\answer{2}
